% Section 5 -- Results: which models can do noi lai, and what predicts failure.
% Budget: 1.50 pages, holds Table 2 and Figures 2-3.
% STATUS: hypotheses only, in the wording of docs/DESIGN_DECISIONS.md section 1. No model has
% been run. Every result is a \placeholder.
\section{Results: Who Can Do Nói Lái, and What Predicts Failure}
\label{sec:results}

\paragraph{Pre-registered hypotheses.}
\hyp{1} (\emph{alignment beyond count}): within models, alignment has a positive and tokens per syllable a negative coefficient, and adding the split indicator and alignment to a model with tokens per syllable improves fit, decided by the nested likelihood-ratio test of Section~\ref{sec:setup}; a model with fewer than 300 misaligned split syllables in the main sample contributes descriptively only, and if fewer than four tokenizer families clear that floor \hyp{1} is descriptive. \hyp{2} is a descriptive case study, not a tested hypothesis, of PhoGPT-4B-Chat's whole-syllable vocabulary. \hyp{6} (\emph{memorization}): exact attested items outscore generated items matched on variant, task, lexicality, spelling trigger and tone class, and a guided-instruction probe \citep{2308.08493} recovers attested but not matched items; it is tested only if at least 100 native-verified exact two-syllable rows from three collections exist at the data freeze.

\paragraph{Overall accuracy.}
\placeholder{Table~\ref{tab:main} reading: models above the human band and at the wrong-variant floor; variant ordering on non-degenerate items (expected \var{2} easiest, \var{3} hardest); \task{T3} against 50\%.}

\begin{table*}[t]
  \centering\small\setlength{\tabcolsep}{3.5pt}
  \begin{tabular}{@{}lrrrrrrr@{}}
    \toprule
    & \multicolumn{4}{c}{\task{T1} transformation, strict (by variant)} & \task{T2} & \multicolumn{2}{c}{\task{T3} (no spelling twins)} \\
    \cmidrule(lr){2-5}\cmidrule(lr){6-6}\cmidrule(lr){7-8}
    Model & \var{1} & \var{2} & \var{3} & \var{4} & decoding & balanced acc. & log-prob.\ (yes/no; pair) \\
    \midrule
    Mean human (\placeholder{$n$} raters, 246 items) & \placeholder{--} & \placeholder{--} & \placeholder{--} & \placeholder{--} & \placeholder{--} & \placeholder{--} & -- \\
    Copy / reversal & 0.0 & 0.0 & 0.0 & 0.0 & 0.0 & 50.0 & 50.0 \\
    Wrong-variant rule & \placeholder{--} & \placeholder{--} & \placeholder{--} & \placeholder{--} & \placeholder{--} & \placeholder{--} & -- \\
    \midrule
    \placeholder{per model} & & & & & & & \\
    \bottomrule
  \end{tabular}
  \caption{Accuracy (\%) with 95\% base-pair bootstrap intervals, mean over three paraphrases (range, lenient \task{T1} and spelling twins in Appendix~\ref{app:extra}). Human row: mean rater on the 246 double-judged items, on which every model is also scored. Log-prob.: open models only, never compared with the generated answer. \placeholder{TODO: filled by \texttt{noilai.stats.tables} from the scored runs.}}
  \label{tab:main}
\end{table*}

\paragraph{What predicts failure (E2).}
\placeholder{Figure~\ref{fig:alignment} reading; verdict on \hyp{1} (coefficients, likelihood-ratio test, floor counts, GEE agreement, share of per-model slopes with the predicted sign); the tokenizer-fixed SEA-LION pair (suggests, never isolates); \hyp{2}.}

\begin{figure}[t]
  \centering
  \fbox{\parbox{0.95\columnwidth}{\centering\vspace{1.8cm}\placeholder{Figure 2: accuracy vs.\ boundary alignment and tokens per syllable; coefficients with bootstrap intervals}\vspace{1.8cm}}}
  \caption{\placeholder{Item-level accuracy as a function of boundary alignment among split syllables (left) and tokens per syllable (right), pooled over models with 95\% intervals; coefficients from the fixed-effect fit with base-pair cluster-bootstrap intervals (inset).}}
  \label{fig:alignment}
\end{figure}

\paragraph{Explicit input.}
\placeholder{Decomposed vs.\ raw input per model: near ceiling places the bottleneck on access, low on the rule.}

\paragraph{Error taxonomy.}
\placeholder{Figure~\ref{fig:errors} reading against the registered expectations: spelling errors on trigger items, copy and reversal in the smallest models.}

\begin{figure}[t]
  \centering
  \fbox{\parbox{0.95\columnwidth}{\centering\vspace{1.8cm}\placeholder{Figure 3: error breakdown by model group}\vspace{1.8cm}}}
  \caption{\placeholder{Error classes as shares of all \task{T1} outputs, by model group and variant.}}
  \label{fig:errors}
\end{figure}

\paragraph{Ablations and memorization.}
\placeholder{Ablations; \hyp{6} with balance and the four discriminators.}
